% Curs 10 - Pagina 4 (6-XII-1994).
% Fragment pentru % Curs 10 - Pagina 4 (6-XII-1994).
% Fragment pentru % Curs 10 - Pagina 4 (6-XII-1994).
% Fragment pentru % Curs 10 - Pagina 4 (6-XII-1994).
% Fragment pentru \input{courses/pages/curs10-pagina4-ro}.
% Continuarea ecuatiei claselor de la pagina 3; fara inceput de document sau intrerupere fortata de pagina.

\subsection*{Curs 10 - Pagina 4}

\noindent\underline{Aplicație}

\medskip
\noindent\underline{T. lui Cauchy}: $G$ grup finit și $|G|=n$;
$p\mid |G|$, $p$ prim.

\medskip
\noindent Atunci $(\exists)\,a\in G$ a.î. $\operatorname{ord}(a)=p$.

\bigskip
\noindent\underline{Dem.}:
\begin{align*}
\text{Caz 1:}\quad & p\mid\operatorname{ord}(Z_G),\\
\text{Caz 2:}\quad & p\nmid\operatorname{ord}(Z_G).
\end{align*}

\noindent\underline{Caz 1}: $\Longrightarrow (\exists)\,a\in Z_G$
a.î. $\operatorname{ord}(a)=p$ (deoarece $Z_G$ este grup abelian).

\medskip
\noindent\underline{Caz 2}: $(\exists)\,z_j\in T$, $z_j\notin Z_G$
a.î. $p\nmid [G:C_G(z_j)]$.

\[
\Longrightarrow\quad p\mid\operatorname{ord}(C_G(z_j))
\quad\bigl(\operatorname{ord}(G)
=\operatorname{ord}(C_G(z_j))\,[G:C_G(z_j)]\bigr).
\]

\noindent Dar $\operatorname{ord}(C_G(z_j))<\operatorname{ord}(G)$.

\medskip
\noindent Prin inducție după $n$:
\[
<n\ \longrightarrow\ n.
\]
% Ultima indicatie a inductiei este doar schematica in manuscris.
% Justificarea separata a cazului abelian nu este scrisa aici.

\bigskip
\hrule
\bigskip

\noindent\underline{Lema 1}: Fie $G$ grup $p$, $|G|=p^n$,
$n\geq 1$, și $A$ o $G$-mulțime finită.

\medskip
\noindent Fie
\[
A_0=\{x\in A\mid ax=x,\ (\forall)\,a\in G\}
   =\{x\in A\mid \operatorname{Stab}_G(x)=G\}.
\]

\noindent Atunci
\[
|A|\equiv |A_0|\pmod p.
\]

\bigskip
\noindent\underline{Caz particular}: $A=G$, $|G|=p^n$
și acțiunea prin conjugare.
\[
\Longrightarrow\quad |G|\equiv |Z_G|\pmod p.
\]

\bigskip
\noindent\underline{Dem.}: Fie $T=\{z_1,\ldots,z_q\}$
o transversală pentru relația de echivalență a acțiunii lui $G$ pe $A$.

\[
A_0\subseteq T,\qquad
x\in A_0\ \Longrightarrow\
|\widehat{x}|=[G:\operatorname{Stab}_G(x)]=1
=|\operatorname{Orb}_G(x)|.
\]

\[
|A|=|A_0|+
\sum_{\substack{z_i\in T\\z_i\notin A_0}}
[G:\operatorname{Stab}_G(z_i)].
\]

\noindent Pentru $z_i\notin A_0$:
\[
[G:\operatorname{Stab}_G(z_i)]\mid |G|,
\qquad [G:\operatorname{Stab}_G(z_i)]>1.
\]
\[
\Longrightarrow\quad p\mid[G:\operatorname{Stab}_G(z_i)].
\]

\noindent Deci
\[
p\mid\bigl(|A|-|A_0|\bigr).
\]
.
% Continuarea ecuatiei claselor de la pagina 3; fara inceput de document sau intrerupere fortata de pagina.

\subsection*{Curs 10 - Pagina 4}

\noindent\underline{Aplicație}

\medskip
\noindent\underline{T. lui Cauchy}: $G$ grup finit și $|G|=n$;
$p\mid |G|$, $p$ prim.

\medskip
\noindent Atunci $(\exists)\,a\in G$ a.î. $\operatorname{ord}(a)=p$.

\bigskip
\noindent\underline{Dem.}:
\begin{align*}
\text{Caz 1:}\quad & p\mid\operatorname{ord}(Z_G),\\
\text{Caz 2:}\quad & p\nmid\operatorname{ord}(Z_G).
\end{align*}

\noindent\underline{Caz 1}: $\Longrightarrow (\exists)\,a\in Z_G$
a.î. $\operatorname{ord}(a)=p$ (deoarece $Z_G$ este grup abelian).

\medskip
\noindent\underline{Caz 2}: $(\exists)\,z_j\in T$, $z_j\notin Z_G$
a.î. $p\nmid [G:C_G(z_j)]$.

\[
\Longrightarrow\quad p\mid\operatorname{ord}(C_G(z_j))
\quad\bigl(\operatorname{ord}(G)
=\operatorname{ord}(C_G(z_j))\,[G:C_G(z_j)]\bigr).
\]

\noindent Dar $\operatorname{ord}(C_G(z_j))<\operatorname{ord}(G)$.

\medskip
\noindent Prin inducție după $n$:
\[
<n\ \longrightarrow\ n.
\]
% Ultima indicatie a inductiei este doar schematica in manuscris.
% Justificarea separata a cazului abelian nu este scrisa aici.

\bigskip
\hrule
\bigskip

\noindent\underline{Lema 1}: Fie $G$ grup $p$, $|G|=p^n$,
$n\geq 1$, și $A$ o $G$-mulțime finită.

\medskip
\noindent Fie
\[
A_0=\{x\in A\mid ax=x,\ (\forall)\,a\in G\}
   =\{x\in A\mid \operatorname{Stab}_G(x)=G\}.
\]

\noindent Atunci
\[
|A|\equiv |A_0|\pmod p.
\]

\bigskip
\noindent\underline{Caz particular}: $A=G$, $|G|=p^n$
și acțiunea prin conjugare.
\[
\Longrightarrow\quad |G|\equiv |Z_G|\pmod p.
\]

\bigskip
\noindent\underline{Dem.}: Fie $T=\{z_1,\ldots,z_q\}$
o transversală pentru relația de echivalență a acțiunii lui $G$ pe $A$.

\[
A_0\subseteq T,\qquad
x\in A_0\ \Longrightarrow\
|\widehat{x}|=[G:\operatorname{Stab}_G(x)]=1
=|\operatorname{Orb}_G(x)|.
\]

\[
|A|=|A_0|+
\sum_{\substack{z_i\in T\\z_i\notin A_0}}
[G:\operatorname{Stab}_G(z_i)].
\]

\noindent Pentru $z_i\notin A_0$:
\[
[G:\operatorname{Stab}_G(z_i)]\mid |G|,
\qquad [G:\operatorname{Stab}_G(z_i)]>1.
\]
\[
\Longrightarrow\quad p\mid[G:\operatorname{Stab}_G(z_i)].
\]

\noindent Deci
\[
p\mid\bigl(|A|-|A_0|\bigr).
\]
.
% Continuarea ecuatiei claselor de la pagina 3; fara inceput de document sau intrerupere fortata de pagina.

\subsection*{Curs 10 - Pagina 4}

\noindent\underline{Aplicație}

\medskip
\noindent\underline{T. lui Cauchy}: $G$ grup finit și $|G|=n$;
$p\mid |G|$, $p$ prim.

\medskip
\noindent Atunci $(\exists)\,a\in G$ a.î. $\operatorname{ord}(a)=p$.

\bigskip
\noindent\underline{Dem.}:
\begin{align*}
\text{Caz 1:}\quad & p\mid\operatorname{ord}(Z_G),\\
\text{Caz 2:}\quad & p\nmid\operatorname{ord}(Z_G).
\end{align*}

\noindent\underline{Caz 1}: $\Longrightarrow (\exists)\,a\in Z_G$
a.î. $\operatorname{ord}(a)=p$ (deoarece $Z_G$ este grup abelian).

\medskip
\noindent\underline{Caz 2}: $(\exists)\,z_j\in T$, $z_j\notin Z_G$
a.î. $p\nmid [G:C_G(z_j)]$.

\[
\Longrightarrow\quad p\mid\operatorname{ord}(C_G(z_j))
\quad\bigl(\operatorname{ord}(G)
=\operatorname{ord}(C_G(z_j))\,[G:C_G(z_j)]\bigr).
\]

\noindent Dar $\operatorname{ord}(C_G(z_j))<\operatorname{ord}(G)$.

\medskip
\noindent Prin inducție după $n$:
\[
<n\ \longrightarrow\ n.
\]
% Ultima indicatie a inductiei este doar schematica in manuscris.
% Justificarea separata a cazului abelian nu este scrisa aici.

\bigskip
\hrule
\bigskip

\noindent\underline{Lema 1}: Fie $G$ grup $p$, $|G|=p^n$,
$n\geq 1$, și $A$ o $G$-mulțime finită.

\medskip
\noindent Fie
\[
A_0=\{x\in A\mid ax=x,\ (\forall)\,a\in G\}
   =\{x\in A\mid \operatorname{Stab}_G(x)=G\}.
\]

\noindent Atunci
\[
|A|\equiv |A_0|\pmod p.
\]

\bigskip
\noindent\underline{Caz particular}: $A=G$, $|G|=p^n$
și acțiunea prin conjugare.
\[
\Longrightarrow\quad |G|\equiv |Z_G|\pmod p.
\]

\bigskip
\noindent\underline{Dem.}: Fie $T=\{z_1,\ldots,z_q\}$
o transversală pentru relația de echivalență a acțiunii lui $G$ pe $A$.

\[
A_0\subseteq T,\qquad
x\in A_0\ \Longrightarrow\
|\widehat{x}|=[G:\operatorname{Stab}_G(x)]=1
=|\operatorname{Orb}_G(x)|.
\]

\[
|A|=|A_0|+
\sum_{\substack{z_i\in T\\z_i\notin A_0}}
[G:\operatorname{Stab}_G(z_i)].
\]

\noindent Pentru $z_i\notin A_0$:
\[
[G:\operatorname{Stab}_G(z_i)]\mid |G|,
\qquad [G:\operatorname{Stab}_G(z_i)]>1.
\]
\[
\Longrightarrow\quad p\mid[G:\operatorname{Stab}_G(z_i)].
\]

\noindent Deci
\[
p\mid\bigl(|A|-|A_0|\bigr).
\]
.
% Continuarea ecuatiei claselor de la pagina 3; fara inceput de document sau intrerupere fortata de pagina.

\subsection*{Curs 10 - Pagina 4}

\noindent\underline{Aplicație}

\medskip
\noindent\underline{T. lui Cauchy}: $G$ grup finit și $|G|=n$;
$p\mid |G|$, $p$ prim.

\medskip
\noindent Atunci $(\exists)\,a\in G$ a.î. $\operatorname{ord}(a)=p$.

\bigskip
\noindent\underline{Dem.}:
\begin{align*}
\text{Caz 1:}\quad & p\mid\operatorname{ord}(Z_G),\\
\text{Caz 2:}\quad & p\nmid\operatorname{ord}(Z_G).
\end{align*}

\noindent\underline{Caz 1}: $\Longrightarrow (\exists)\,a\in Z_G$
a.î. $\operatorname{ord}(a)=p$ (deoarece $Z_G$ este grup abelian).

\medskip
\noindent\underline{Caz 2}: $(\exists)\,z_j\in T$, $z_j\notin Z_G$
a.î. $p\nmid [G:C_G(z_j)]$.

\[
\Longrightarrow\quad p\mid\operatorname{ord}(C_G(z_j))
\quad\bigl(\operatorname{ord}(G)
=\operatorname{ord}(C_G(z_j))\,[G:C_G(z_j)]\bigr).
\]

\noindent Dar $\operatorname{ord}(C_G(z_j))<\operatorname{ord}(G)$.

\medskip
\noindent Prin inducție după $n$:
\[
<n\ \longrightarrow\ n.
\]
% Ultima indicatie a inductiei este doar schematica in manuscris.
% Justificarea separata a cazului abelian nu este scrisa aici.

\bigskip
\hrule
\bigskip

\noindent\underline{Lema 1}: Fie $G$ grup $p$, $|G|=p^n$,
$n\geq 1$, și $A$ o $G$-mulțime finită.

\medskip
\noindent Fie
\[
A_0=\{x\in A\mid ax=x,\ (\forall)\,a\in G\}
   =\{x\in A\mid \operatorname{Stab}_G(x)=G\}.
\]

\noindent Atunci
\[
|A|\equiv |A_0|\pmod p.
\]

\bigskip
\noindent\underline{Caz particular}: $A=G$, $|G|=p^n$
și acțiunea prin conjugare.
\[
\Longrightarrow\quad |G|\equiv |Z_G|\pmod p.
\]

\bigskip
\noindent\underline{Dem.}: Fie $T=\{z_1,\ldots,z_q\}$
o transversală pentru relația de echivalență a acțiunii lui $G$ pe $A$.

\[
A_0\subseteq T,\qquad
x\in A_0\ \Longrightarrow\
|\widehat{x}|=[G:\operatorname{Stab}_G(x)]=1
=|\operatorname{Orb}_G(x)|.
\]

\[
|A|=|A_0|+
\sum_{\substack{z_i\in T\\z_i\notin A_0}}
[G:\operatorname{Stab}_G(z_i)].
\]

\noindent Pentru $z_i\notin A_0$:
\[
[G:\operatorname{Stab}_G(z_i)]\mid |G|,
\qquad [G:\operatorname{Stab}_G(z_i)]>1.
\]
\[
\Longrightarrow\quad p\mid[G:\operatorname{Stab}_G(z_i)].
\]

\noindent Deci
\[
p\mid\bigl(|A|-|A_0|\bigr).
\]
